\section{¿Por qué este curso?}

\begin{frame}
	\frametitle{\insertsection}

	Muchos fenómenos en la naturaleza se describen mediante ecuaciones
	diferenciales parciales (EDPs).

	\begin{columns}
		\begin{column}{.5\paperwidth}
			\begin{block}{Los fenómenos físicos se modelan con las EDPs}
				\begin{description}
					\item[Propagación de ondas:]

					      tsunamis, terremotos, sonido.

					\item[Difusión térmica:]

					      distribución de calor en materiales.

					\item[Dinámica de fluidos:]

					      océanos, atmósfera, aire en tuberías.

					\item[Reacciones químicas:]

					      combustión, ecosistemas.
				\end{description}
			\end{block}

			La mayoría de los problema de valor inicial / frontera no
			tienen soluciones analíticas.

			\begin{center}\large
				¡Necesitamos de los métodos numéricos para resolverlas!
			\end{center}
		\end{column}
		\begin{column}{.4\paperwidth}
			\begin{figure}[ht!]
				\centering
				\includegraphics[width=0.4\paperwidth]{geoclaw}
				\caption{Simulación numérica de un tsunami (Fuente: \url{https://dlgeorge.github.io/project/geoclaw-project})}
				\label{fig:tsunami}
			\end{figure}
		\end{column}
	\end{columns}
\end{frame}
